\chapter{The Fundamental Interactions}\label{ch:g11:fundamental-interactions}

Rub a balloon on your hair and it picks up scraps of paper --- beating,
with a few square centimetres of rubber, the gravitational pull of the
entire planet. That easy victory is why atoms hold together and why
heavy nuclei eventually shatter (\cref{ch:g11:nucleus-radioactivity}).
The universe runs on exactly four interactions; this chapter meets
them, weighs them against each other, and assigns each its floor.

\section{Matter from quarks to galaxies}

\begin{definition}[The structure ladder]\label{def:g11:fundamental-interactions:ladder}
Matter is built in floors, each assembled from the one below: three
\emph{quarks}\index{quark} (pointlike, smaller than \qty{e-18}{m}) bind
into a \emph{nucleon}\index{nucleon} --- proton or neutron, about
\qty{e-15}{m} across; nucleons pack into a \emph{nucleus}\index{nucleus}
(\qty{e-15}{m} to \qty{e-14}{m}); a nucleus plus its electron cloud is
an \emph{atom}\index{atom} (\qty{e-10}{m}); atoms bind into
\emph{molecules}\index{molecule} and crystals (\qty{e-9}{m} up), which
assemble into everyday objects (\qty{1}{m}), planets (\qty{e7}{m}),
planetary systems (\qty{e11}{m}) and galaxies (\qty{e21}{m}).
\end{definition}

\begin{definition}[The four interactions]\label{def:g11:fundamental-interactions:four}
A \emph{fundamental interaction}\index{fundamental interaction} is a
force not reducible to other forces; there are exactly four:
\emph{gravitation}, the \emph{electromagnetic interaction}\index{electromagnetic interaction},
the strong interaction and the weak interaction.
\end{definition}

\begin{remark}\label{rem:g11:fundamental-interactions:empty}
The floors are separated by voids: an atom is $10^{5}$ times wider than
its nucleus --- scale the nucleus up to a \qty{1}{cm} marble and the
electron cloud is a kilometre across. Something must hold these sparse
floors together, and gravity is almost never it.
\end{remark}

\begin{omfigure}
\begin{tikzpicture}[font=\small, x=0.26cm]
  \draw[-Stealth] (-1,0) -- (41.5,0);
  \foreach \x/\l/\d in {0/{10^{-18}}/2pt, 3/{10^{-15}}/13pt,
                        8/{10^{-10}}/2pt, 18/{10^{0}}/2pt,
                        25/{10^{7}}/2pt, 29/{10^{11}}/13pt,
                        39/{10^{21}}/2pt}
    \draw (\x,0.1) -- (\x,-0.1) node[below=\d] {$\l$};
  \foreach \x/\n in {0/quark, 3/nucleus, 8/atom, 18/human, 25/Earth,
                     29/{Sun--Earth}, 39/galaxy}{
    \fill[omDef] (\x,0) circle (1.6pt);
    \node[rotate=45, anchor=west, inner sep=1.5pt] at (\x,0.24) {\n};}
\end{tikzpicture}

{\small The ladder of structure on a powers-of-ten axis: thirty-nine
decades separate the quark from a galaxy. Sizes in metres.}
\end{omfigure}

\section{Electric charge}

\begin{definition}[Electric charge]\label{def:g11:fundamental-interactions:charge}
\emph{Electric charge}\index{electric charge} is the property of matter
on which the electromagnetic interaction acts, as mass is the property
gravitation acts on. It is measured in \emph{coulombs}\index{coulomb}
(\unit{C}) and comes in \emph{two signs}, positive and negative, which
cancel when added; like charges repel, unlike charges attract.
\end{definition}

\begin{definition}[Elementary charge]\label{def:g11:fundamental-interactions:elementary}
Charge is grained: every observed charge is a whole multiple of the
\emph{elementary charge}\index{elementary charge}
$e = \qty{1.60e-19}{C}$. The proton carries $+e$, the electron $-e$, the
neutron $0$; a charge $q$ holds $N = \abs{q}/e$ elementary charges.
\end{definition}

\begin{proposition}[Conservation of charge]\label{prop:g11:fundamental-interactions:conservation}
The total charge of an isolated system never changes: charge is neither
created nor destroyed, only transferred --- in practice by electrons,
the lightest carriers.
\end{proposition}
\admitted

\begin{remark}\label{rem:g11:fundamental-interactions:conservation}
No violation has ever been observed; the deep reason, a symmetry of
electromagnetism, is given in the university volumes.
\end{remark}

\begin{definition}[Charging by friction and by contact]\label{def:g11:fundamental-interactions:charging}
Rubbing two insulators tears electrons off one and onto the other,
leaving opposite charges of equal magnitude (\emph{charging by
friction}\index{charging by friction}); touching a charged object to a
neutral one shares the charge (\emph{charging by
contact}\index{charging by contact}). Only electrons move; the total
charge is conserved.
\end{definition}

\begin{example}[Counting electrons]\label{ex:g11:fundamental-interactions:balloon}
A balloon rubbed on hair acquires $q = \qty{-1.6e-8}{C}$: it has
\emph{gained} $N = \num{1.6e-8} / \num{1.6e-19} = \num{1.0e11}$
electrons, and the hair is left with exactly $+\qty{1.6e-8}{C}$.
\end{example}

\section{Coulomb's law}

\begin{theorem}[Coulomb's law]\label{thm:g11:fundamental-interactions:coulomb}
Two point charges $q_1$ and $q_2$ a distance $d$ apart exert on each
other forces along the line joining them, of equal magnitudes and
opposite directions --- repulsive for like signs, attractive for unlike
signs --- with
\[
F = k\,\frac{\abs{q_1 q_2}}{d^{2}},
\qquad k = \qty{8.99e9}{N.m^2/C^2}.
\]
\end{theorem}
\admitted

\begin{remark}\label{rem:g11:fundamental-interactions:torsion}
Coulomb established the law in 1785 by measuring the twist of a thin
torsion fibre; why the exponent is exactly $2$ is answered in the university volumes.
\end{remark}

\begin{omfigure}
\begin{tikzpicture}[font=\small]
  \draw[omDef, thick, fill=omDef!20] (0,0) circle (0.3) node[black] {$+$};
  \draw[omDef, thick, fill=omDef!20] (2.6,0) circle (0.3) node[black] {$+$};
  \draw[-Stealth, omThm, very thick] (-0.3,0) -- (-1.5,0)
    node[midway, above] {$\vect F_{2/1}$};
  \draw[-Stealth, omThm, very thick] (2.9,0) -- (4.1,0)
    node[midway, above] {$\vect F_{1/2}$};
\end{tikzpicture}\qquad
\begin{tikzpicture}[font=\small]
  \draw[omDef, thick, fill=omDef!20] (0,0) circle (0.3) node[black] {$+$};
  \draw[omProp, thick, fill=omProp!25] (2.6,0) circle (0.3) node[black] {$-$};
  \draw[-Stealth, omThm, very thick] (0.3,0) -- (1.1,0)
    node[midway, above] {$\vect F_{2/1}$};
  \draw[-Stealth, omThm, very thick] (2.3,0) -- (1.5,0)
    node[midway, above] {$\vect F_{1/2}$};
\end{tikzpicture}

{\small Like charges repel (left), unlike charges attract (right); the
two forces always have equal magnitudes and opposite directions.}
\end{omfigure}

\begin{remark}[A formal twin of gravitation]\label{rem:g11:fundamental-interactions:twin}
Coulomb's law is, symbol for symbol, the law of universal gravitation
(\cref{ch:g10:universal-gravitation}) with charges in place of masses:
\begin{center}
\begin{tabular}{lll}
source & mass $m$ & charge $q$ \\
law & $F = G\,\dfrac{m_1 m_2}{d^2}$ & $F = k\,\dfrac{\abs{q_1 q_2}}{d^2}$ \\[1.5ex]
constant & $G = \qty{6.67e-11}{N.m^2/kg^2}$ & $k = \qty{8.99e9}{N.m^2/C^2}$ \\
sign & always attractive & attractive or repulsive \\
\end{tabular}
\end{center}
The two differences drive everything below: electricity is
overwhelmingly stronger, and it alone can cancel itself.
\end{remark}

\begin{example}[Two protons, both forces at once]\label{ex:g11:fundamental-interactions:protons}
Two protons ($m_p = \qty{1.67e-27}{kg}$) sit $d = \qty{1.0e-15}{m}$
apart --- nuclear neighbours. Electric repulsion:
\[
F_E = \num{8.99e9} \times
\frac{(\num{1.60e-19})^{2}}{(\num{1.0e-15})^{2}}
\approx \qty{2.3e2}{N}
\]
--- the weight of a $23$-kg suitcase, carried by a particle of
$10^{-27}$ kilograms. Gravitational attraction:
$F_G = \num{6.67e-11} \times (\num{1.67e-27})^{2} /
(\num{1.0e-15})^{2} \approx \qty{1.9e-34}{N}$. The ratio
$F_E / F_G = k e^{2} / (G m_p^{2}) \approx \num{1.2e36}$ is
distance-independent ($d^{2}$ cancels): between elementary particles,
gravity is $10^{36}$ times too weak to matter --- at \emph{any} distance.
\end{example}

\section{The strong and weak interactions}

The example leaves a scandal: a helium nucleus holds two protons that
repel with tens of newtons --- and yet helium exists.

\begin{definition}[The strong interaction]\label{def:g11:fundamental-interactions:strong}
The \emph{strong interaction}\index{strong interaction} binds quarks
into nucleons and nucleons into nuclei. At \qty{e-15}{m} it is
attractive and far stronger than the electric repulsion (thousands of
newtons between nucleons); beyond a few $10^{-15}$ metres it vanishes:
its \emph{range}\index{range (of an interaction)} is about
\qty{e-15}{m}. It grips protons and neutrons alike and ignores charge.
\end{definition}

\begin{omfigure}
\begin{tikzpicture}[font=\small]
  \draw[omDef, thick, fill=omDef!20] (0,0) circle (0.34) node[black] {p};
  \draw[omDef, thick, fill=omDef!20] (3.2,0) circle (0.34) node[black] {p};
  \draw[-Stealth, omThm, thick] (-0.34,0) -- (-1.7,0)
    node[midway, above] {$F_E \approx \qty{230}{N}$};
  \draw[-Stealth, omThm, thick] (3.54,0) -- (4.9,0) node[midway, above] {$F_E$};
  \draw[-Stealth, omMeth, very thick] (0.2,0.55) -- (1.45,0.55)
    node[midway, above] {strong};
  \draw[-Stealth, omMeth, very thick] (3.0,0.55) -- (1.75,0.55)
    node[midway, above] {strong};
  \draw[Stealth-Stealth] (0,-0.65) -- (3.2,-0.65)
    node[midway, below] {$d = \qty{1.0e-15}{m}$};
\end{tikzpicture}

{\small The tug-of-war inside a nucleus: at $10^{-15}$~m the strong
attraction (thousands of newtons) beats the electric repulsion.}
\end{omfigure}

\begin{remark}[Why nuclei exist --- and why heavy ones give way]\label{rem:g11:fundamental-interactions:nuclei}
The two forces in the nucleus scale differently. The strong glue is
short-ranged: a nucleon binds only its few neighbours within
\qty{e-15}{m}. The electric repulsion is long-ranged: every proton
pushes on \emph{every other} proton. As nuclei grow, repulsion
accumulates faster than glue: heavy nuclei need extra neutrons (glue
without charge), and beyond about $80$ protons even that fails --- the
heaviest nuclei live on borrowed time (\cref{ch:g11:nucleus-radioactivity}).
\end{remark}

\begin{definition}[The weak interaction]\label{def:g11:fundamental-interactions:weak}
The \emph{weak interaction}\index{weak interaction}, of range shorter
still, lets a neutron turn into a proton (and back): it drives the
$\beta$ radioactivity of \cref{ch:g11:nucleus-radioactivity}.
\end{definition}

\section{Who rules which floor}

\begin{method}[Auditing a scale]\label{met:g11:fundamental-interactions:audit}
To find which interaction rules a structure of size $d$, check range
and cancellation:
\begin{enumerate}
  \item $d \leq \qty{e-14}{m}$: the \emph{strong interaction} rules ---
        it beats the electric repulsion, and gravity is $10^{36}$ out.
  \item atoms to mountains ($\qty{e-10}{m}$ to $\qty{e4}{m}$): the
        strong force is out of range; the \emph{electromagnetic
        interaction} rules --- bonds, cohesion, friction, contact.
  \item planets and beyond ($d \geq \qty{e6}{m}$): matter is neutral to
        fantastic precision, so electricity cancels itself; mass has
        one sign and always adds --- \emph{gravitation} rules the sky.
\end{enumerate}
The weak interaction holds no floor: it binds nothing, but arbitrates
transformations ($\beta$ decay).
\end{method}

\begin{omfigure}
\begin{tikzpicture}[font=\small, x=0.26cm]
  \draw[-Stealth] (-1,0) -- (41.5,0);
  \foreach \x/\l in {3/{10^{-15}}, 8/{10^{-10}}, 18/{10^{0}},
                     24/{10^{6}}, 39/{10^{21}}}
    \draw (\x,0.1) -- (\x,-0.1) node[below] {$\l$};
  \draw[omProp, fill=omProp!35] (3,0.3) rectangle (4,0.75);
  \node[anchor=west, omProp] at (4.4,0.52) {strong};
  \draw[omDef, fill=omDef!25] (8,0.95) rectangle (24,1.4);
  \node[omDef] at (16,1.17) {electromagnetic};
  \draw[omThm, fill=omThm!25] (24,1.6) rectangle (39,2.05);
  \node[omThm] at (31.5,1.82) {gravitation};
\end{tikzpicture}

{\small Who rules where: strong for nuclei, electromagnetic from atoms
to mountains, gravitation for planets and beyond. Sizes in metres.}
\end{omfigure}

\section{Exercises}

\begin{exercise}[$\star$]\label{exo:g11:fundamental-interactions:1}
Give the order of magnitude of each size and its floor of the ladder:
proton \qty{1.7e-15}{m}; water molecule \qty{2.8e-10}{m}; grain of sand
\qty{2e-4}{m}; Earth \qty{1.3e7}{m}; Milky Way \qty{9.5e20}{m}.
\end{exercise}

\begin{exercise}[$\star$]\label{exo:g11:fundamental-interactions:2}
A balloon rubbed on hair carries $q = \qty{-4.8e-9}{C}$. Has it gained
or lost electrons, and how many? What is the hair's charge afterwards,
and which law says so?
\end{exercise}

\begin{exercise}[$\star$]\label{exo:g11:fundamental-interactions:3}
Charges $q_1 = \qty{+2.0e-6}{C}$ and $q_2 = \qty{-3.0e-6}{C}$ sit
\qty{30}{cm} apart. Compute the force; attractive or repulsive?
Compare the forces on $q_1$ and on $q_2$.
\end{exercise}

\begin{exercise}[$\star$]\label{exo:g11:fundamental-interactions:4}
Two charges attract with a force $F_0$. What does it become if: the
distance doubles; the distance is divided by $3$; one charge doubles;
both charges \emph{and} the distance double?
\end{exercise}

\begin{exercise}[$\star$]\label{exo:g11:fundamental-interactions:5}
Name the interaction responsible for: (a) the Moon orbiting the Earth;
(b) the cohesion of a salt crystal; (c) the cohesion of a helium
nucleus; (d) the $\beta$ decay of carbon-14; (e) the balloon of
\cref{exo:g11:fundamental-interactions:2} sticking to a wall.
\end{exercise}

\begin{exercise}[$\star\star$]\label{exo:g11:fundamental-interactions:6}
The two protons of a helium nucleus sit $d = \qty{2.0e-15}{m}$ apart
($m_p = \qty{1.67e-27}{kg}$). Compute their electric repulsion, their
gravitational attraction, and the ratio. What keeps the nucleus whole?
\end{exercise}

\begin{exercise}[$\star\star$]\label{exo:g11:fundamental-interactions:7}
In a hydrogen atom, the electron ($m_e = \qty{9.11e-31}{kg}$) sits
$d = \qty{5.3e-11}{m}$ from the proton ($m_p = \qty{1.67e-27}{kg}$).
Compute the electric and gravitational forces between them and their
ratio. What holds atoms together?
\end{exercise}

\begin{exercise}[$\star\star$]\label{exo:g11:fundamental-interactions:8}
Sphere A carries $q_A = \qty{+8.0e-9}{C}$; an identical sphere B is
neutral. They are touched together, then separated to \qty{10}{cm}.
What charge does each carry (check the conservation of charge)? Compute
the force between them --- attractive or repulsive?
\end{exercise}

\begin{exercise}[$\star\star$]\label{exo:g11:fundamental-interactions:9}
Two identical charges $q = \qty{1.0e-6}{C}$ repel with \qty{0.90}{N}:
how far apart are they? Where does the force drop to \qty{0.10}{N}?
\end{exercise}

\begin{exercise}[$\star\star$]\label{exo:g11:fundamental-interactions:10}
Two protons in neighbouring molecules sit $d = \qty{1.0e-10}{m}$ apart.
Compute their electric repulsion. What does the strong interaction
contribute at this distance? Which interaction binds molecules?
\end{exercise}

\begin{exercise}[$\star\star$]\label{exo:g11:fundamental-interactions:11}
Each of two coins holds about $\num{e23}$ electrons. Move one electron
in a \emph{million} between them and set the coins \qty{1.0}{m} apart.
Compute each coin's charge, then the force. The weight of what mass
equals it? Conclude on the neutrality of everyday matter.
\end{exercise}

\begin{exercise}[$\star\star\star$]\label{exo:g11:fundamental-interactions:12}
A uranium nucleus has diameter \qty{1.4e-14}{m} and $92$ protons.
Compute the repulsion between two protons at opposite ends. Can the
strong interaction bind this pair --- why not? How many repelling
proton pairs are there? Why do heavy nuclei need extra neutrons, and
why beyond some size can even neutrons not save them?
\end{exercise}

\begin{exercise}[$\star\star\star$]\label{exo:g11:fundamental-interactions:13}
Two balls of mass \qty{1.0}{g} hang from \qty{50}{cm} threads tied to
the same point. Given the same charge $q$, they settle \qty{6.0}{cm}
apart. Each ball is in equilibrium under its weight, the thread's
tension and the electric force: show that $F = mg\tan\theta$ ($\theta$:
angle of thread to vertical) and that here $\tan\theta \approx 0.060$.
Compute $F$, then $q$, then the number of elementary charges moved.
\end{exercise}

\begin{exercise}[$\star\star\star$]\label{exo:g11:fundamental-interactions:14}
Data: $M_E = \qty{5.97e24}{kg}$, $M_M = \qty{7.35e22}{kg}$, Earth--Moon
distance $d = \qty{3.84e8}{m}$. Compute the Earth--Moon gravitational
force. Gravity is switched off: what charges $+Q$ on the Earth and $-Q$
on the Moon keep the same attraction? How many electrons is $Q$, and
what do they weigh ($m_e = \qty{9.11e-31}{kg}$)? Comment.
\end{exercise}

\begin{exercise}[$\star\star\star$]\label{exo:g11:fundamental-interactions:15}
In a salt crystal, Na$^{+}$ and Cl$^{-}$ ions (charges $\pm e$) sit
$d = \qty{2.8e-10}{m}$ apart; a Cl$^{-}$ ion has mass \qty{5.9e-26}{kg}.
Compute the force between neighbouring ions and the weight of a
Cl$^{-}$ ion; give the ratio. Then explain how gravity wins at large
scales --- no mountain tops \qty{e4}{m} --- though each bond beats it
by fifteen orders of magnitude.
\end{exercise}

\section{Problem: The audit of the four forces}

\begin{problem}[{Weekend problem --- why the world neither collapses
nor flies apart: the four interactions audited floor by floor, down to
matter's neutrality to two parts in $10^{18}$}]
\label{pb:g11:fundamental-interactions:1}
Atoms do not implode, nuclei mostly hold, the Moon neither crashes nor
escapes; this problem audits who deserves the credit. Data:
$e = \qty{1.60e-19}{C}$,
$k = \qty{8.99e9}{N.m^2/C^2}$, $G = \qty{6.67e-11}{N.m^2/kg^2}$,
$m_p = \qty{1.67e-27}{kg}$, $m_e = \qty{9.11e-31}{kg}$,
$M_E = \qty{5.97e24}{kg}$, $M_M = \qty{7.35e22}{kg}$, Earth--Moon
distance \qty{3.84e8}{m}.

\medskip
\noindent\textbf{Part I --- The ladder.}
\begin{enumerate}
  \item List the floors of the ladder from quark to galaxy, one order
        of magnitude of size each.
  \item Scale model: the nucleus becomes a \qty{1}{cm} marble; how wide
        is the atom?
  \item What fraction of an atom's volume does its nucleus occupy?
  \item A hair is \qty{70}{\micro m} thick. How many atoms span it?
  \item How many powers of ten from the quark ceiling (\qty{e-18}{m})
        to a galaxy (\qty{e21}{m})?
\end{enumerate}

\medskip
\noindent\textbf{Part II --- The electric floor.}
\begin{enumerate}[resume]
  \item In a hydrogen atom the electron sits $d = \qty{5.3e-11}{m}$
        from the proton: compute the electric force between them.
  \item Compute the gravitational force between them, and the ratio of
        the two. Which one holds the atom together?
  \item What changes if gravity is switched off inside atoms? And if
        electricity is?
  \item Why are the ``contact'' forces of everyday life --- the floor
        pushing on your feet, friction --- electromagnetic in disguise?
  \item Each of two coins holds about $\num{e23}$ electrons. Transfer
        one electron in a \emph{billion} between them, set them
        \qty{1.0}{m} apart: compute the charges and the force. Conclusion?
  \item Which interaction rules the floors from \qty{e-10}{m} to
        \qty{e4}{m}, and why does the strong interaction not compete?
\end{enumerate}

\medskip
\noindent\textbf{Part III --- The nuclear floor.}
\begin{enumerate}[resume]
  \item Two protons sit \qty{1.0e-15}{m} apart in a nucleus: compute
        their electric repulsion.
  \item Unrestrained, what acceleration would this force give a proton?
        Compare with $g$.
  \item List the three properties the strong interaction needs to
        explain why nuclei exist, why neutrons are gripped too, and why
        nuclei are tiny.
  \item In a uranium nucleus (diameter \qty{1.4e-14}{m}, $92$ protons),
        compute the repulsion between two protons at opposite ends. Can
        the strong interaction bind that pair directly? Explain why
        repulsion wins as nuclei grow, and what the neutrons are for.
  \item In one sentence, as promised: what does the weak interaction do?
\end{enumerate}

\medskip
\noindent\textbf{Part IV --- The astronomical floor, and the verdict.}
\begin{enumerate}[resume]
  \item Compute the gravitational force between the Earth and the Moon.
  \item The Earth contains about $\num{1.8e51}$ protons (and as many
        electrons), the Moon about $\num{2.2e49}$. What fraction $f$ of
        each body's proton charge, left uncancelled, would make the
        electric force match the gravitational one?
  \item Why does gravitation, the weakling of question 7, rule
        astronomy? Two reasons.
  \item Verdict --- one sentence per floor (nucleus, atom to mountain,
        planet and up): what holds it together? Then answer the title:
        why does the world neither collapse nor fly apart?
\end{enumerate}
\end{problem}
